\documentclass[11pt,a4paper]{article}

% ---------------------------------------------------------------------------
% Packages
% ---------------------------------------------------------------------------
\usepackage[utf8]{inputenc}
\usepackage[T1]{fontenc}
\usepackage[margin=1in]{geometry}
\usepackage{amsmath, amssymb, amsthm}
\usepackage{mathtools}
\usepackage{hyperref}
\usepackage{xcolor}
\usepackage{listings}
\usepackage{booktabs}
\usepackage{enumitem}
\usepackage{graphicx}
\usepackage{caption}
\usepackage{float}
\usepackage[protrusion=true,expansion=false]{microtype}

% ---------------------------------------------------------------------------
% Colours and code listings
% ---------------------------------------------------------------------------
\definecolor{codebg}{rgb}{0.97,0.97,0.97}
\definecolor{codekw}{rgb}{0.12,0.32,0.62}
\definecolor{codestr}{rgb}{0.16,0.50,0.16}
\definecolor{codecomm}{rgb}{0.45,0.45,0.45}
\definecolor{codenum}{rgb}{0.55,0.55,0.55}

\lstdefinestyle{python}{
  language=Python,
  basicstyle=\ttfamily\footnotesize,
  backgroundcolor=\color{codebg},
  keywordstyle=\color{codekw}\bfseries,
  stringstyle=\color{codestr},
  commentstyle=\color{codecomm}\itshape,
  numberstyle=\tiny\color{codenum},
  numbers=left,
  numbersep=8pt,
  frame=single,
  rulecolor=\color{codebg},
  showstringspaces=false,
  breaklines=true,
  columns=fullflexible,
  keepspaces=true,
  xleftmargin=12pt,
  framexleftmargin=10pt,
  framexrightmargin=0pt,
  literate={->}{{$\rightarrow$}}1 {>=}{{$\geq$}}1 {<=}{{$\leq$}}1
}
\lstset{style=python}

% ---------------------------------------------------------------------------
% Theorems
% ---------------------------------------------------------------------------
\newtheorem{definition}{Definition}[section]
\newtheorem{theorem}[definition]{Theorem}
\newtheorem{proposition}[definition]{Proposition}
\newtheorem{lemma}[definition]{Lemma}
\newtheorem{example}[definition]{Example}
\newtheorem{remark}[definition]{Remark}

% ---------------------------------------------------------------------------
% Macros
% ---------------------------------------------------------------------------
\newcommand{\Rplus}{\mathbb{R}_{\geq 0}}
\newcommand{\Rplustop}{\overline{\Rplus}}
\newcommand{\Nat}{\mathbb{N}}
\newcommand{\Min}{\mathsf{Min}}
\newcommand{\Antichain}{\mathsf{A}}
\newcommand{\topp}{\top}
\newcommand{\codename}{\texttt{codesign-mcdp}}

\hypersetup{
  colorlinks=true,
  linkcolor=blue!60!black,
  citecolor=blue!60!black,
  urlcolor=blue!60!black,
}

% ---------------------------------------------------------------------------
% Title
% ---------------------------------------------------------------------------
\title{\codename: A Python Library for \\ Monotone Co-Design Problems \\
       \vspace{0.4em} \large Version 0.1.0 -- Reference Manual}
\author{}
\date{May 2026}

\begin{document}
\maketitle

\begin{abstract}
\noindent \codename{} is a Python library for formulating and solving
\emph{Monotone Co-Design Problems} (MCDPs) in the framework of Censi
(2015). A design problem is a relation between two posets, a functionality
poset $F$ and a resource poset $R$; given a target functionality, the
problem asks for the antichain of minimal resources needed to deliver it.
Design problems compose under three operators (series, parallel,
feedback), and the resulting class is closed under composition. The
library implements the antichain calculus, six primitive design-problem
types, the three composition operators, a Kleene fixed-point solver, two
high-level builders (an MCDPL-style declarative builder and a modular
\texttt{System} builder), and a collection of worked examples. This
document is the reference manual for version 0.1.0.
\end{abstract}

\tableofcontents
\bigskip

% ===========================================================================
\section{Introduction}
% ===========================================================================

A \emph{co-design problem} is the joint design of components whose
specifications mutually depend on each other. A canonical example: a
battery must store enough energy to power an actuator, but the actuator
must lift the battery itself, so the battery's mass appears in the
energy it has to supply. There is no way to fix one without fixing the
other.

Censi (2015) showed that under modest monotonicity assumptions, such
problems have a clean algebraic theory: the relations between resources
and functionalities form a category closed under series, parallel, and
feedback composition, and solutions are computed by a Kleene
fixed-point iteration in the lattice of antichains. The original
publication included a software tool, \texttt{MCDPL}, distributed
through \texttt{co-design.science}, which exposes the framework via a
domain-specific language.

\codename{} is a from-scratch Python implementation of the algorithmic
core, designed around composable Python objects rather than a separate
DSL. It targets users who want to formulate and solve MCDPs from inside
their existing Python codebase without context switches to a DSL or
external solver. The library is small (about 2000 lines of Python), has
no required runtime dependencies beyond the standard library, and is
licensed under MIT.

\paragraph{Scope of this document.} This manual covers the public API
and the mathematical structure it implements. Section~\ref{sec:theory}
recalls the formal background. Sections~\ref{sec:posets} through
\ref{sec:solver} describe the data types, primitive DPs, composition
operators, and the solver in detail. Sections~\ref{sec:mcdp-builder}
and~\ref{sec:system-builder} document the two high-level builders.
Section~\ref{sec:examples} walks through every example shipped with the
library. Section~\ref{sec:guidelines} collects modelling guidelines and
common pitfalls. Section~\ref{sec:limitations} discusses limitations
and future directions.

\paragraph{Conventions.} Code listings are typeset in monospaced font.
Type signatures use Python's standard notation
($\texttt{f: F} \rightarrow \texttt{R}$, for example). Mathematical
quantities use $F$, $R$, $h$, $\top$, $\bot$, $\Min$, etc.

% ===========================================================================
\section{Mathematical Background}
\label{sec:theory}
% ===========================================================================

This section recalls the formal definitions from Censi (2015) needed
to specify the API. Proofs are omitted; the reader is referred to the
paper for the full development.

\subsection{Posets and Lattices}

\begin{definition}[Poset]
A \emph{poset} (partially ordered set) is a pair $(P, \leq)$ where
$\leq$ is a reflexive, antisymmetric, transitive binary relation on
the set $P$.
\end{definition}

\begin{definition}[Bottom and top]
An element $\bot \in P$ is a \emph{bottom} (least element) if
$\bot \leq x$ for every $x \in P$. An element $\topp \in P$ is a
\emph{top} (greatest element) if $x \leq \topp$ for every $x \in P$.
A poset has at most one bottom and at most one top.
\end{definition}

\begin{example}[Reals and Naturals]
The non-negative reals augmented with infinity, $\Rplustop = \Rplus
\cup \{+\infty\}$, ordered by the usual $\leq$, form a poset with
$\bot = 0$ and $\topp = +\infty$. Likewise $\overline{\Nat} =
\Nat \cup \{+\infty\}$.
\end{example}

\begin{example}[Product order]
Given posets $(P_1, \leq_1), \dots, (P_n, \leq_n)$, the product
$P_1 \times \cdots \times P_n$ with the componentwise order
\[
(x_1, \dots, x_n) \leq (y_1, \dots, y_n)
  \iff \forall i.\, x_i \leq_i y_i
\]
is a poset. The bottom is $(\bot_1, \dots, \bot_n)$ and the top is
$(\topp_1, \dots, \topp_n)$.
\end{example}

\begin{definition}[Antichain]
A subset $A \subseteq P$ is an \emph{antichain} if no two distinct
elements of $A$ are comparable: $\forall x, y \in A,\ x \leq y \implies x = y$.
We write $\Antichain[P]$ for the set of all antichains of $P$.
\end{definition}

\begin{definition}[Min operator]
For any subset $S \subseteq P$, the \emph{Min operator} returns the
set of minimal elements:
\[
\Min(S) = \{ x \in S : \nexists y \in S.\, y < x \}.
\]
$\Min(S)$ is always an antichain. For a finite $S$, $\Min(S)$ is the
Pareto front of $S$.
\end{definition}

\subsection{Antichains as a Lattice}

The set of antichains of a poset carries a natural order: $A \leq A'$
in $\Antichain[P]$ iff every point in $A'$ is dominated by some point
in $A$. Intuitively, $A \leq A'$ means $A$ is "lower" or "easier",
having weaker Pareto-minimal demands.

\begin{proposition}
$(\Antichain[P], \leq)$ is a complete lattice when $P$ is finite or
chain-complete. Joins are computed by union followed by $\Min$. The
bottom of $\Antichain[P]$ is $\{\bot_P\}$ (or $\emptyset$ if $P$ has
no bottom); the top is $\emptyset$.
\end{proposition}

This lattice is the state space of the Kleene iteration (Section
\ref{sec:loop}).

\subsection{Design Problems}

\begin{definition}[Design Problem]
A \emph{design problem (DP)} is a tuple $\langle F, R, h \rangle$
where $F$ and $R$ are posets and $h : F \to \Antichain[R]$ is a
function from functionality requests to antichains of resources.
$h(f)$ is interpreted as the Pareto front of minimal resources
required to deliver functionality $f$.
\end{definition}

\begin{definition}[Monotone Design Problem]
A DP $\langle F, R, h \rangle$ is \emph{monotone} (an MCDP) if $h$
is monotone with respect to the antichain order: $f \leq f'$ in $F$
implies $h(f) \leq h(f')$ in $\Antichain[R]$.
\end{definition}

Intuitively, monotonicity says that demanding more functionality
cannot reduce the minimal resources needed. This corresponds to the
common-sense engineering principle that you don't get more for less.

\subsection{Composition}

The class of MCDPs is closed under three operators.

\begin{definition}[Series]
Given DPs $d_1 = \langle F, M, h_1 \rangle$ and $d_2 = \langle M, R,
h_2 \rangle$ whose interface poset $M$ matches, their \emph{series
composition} is
\[
(d_1 \circ d_2)(f) = \Min\!\left( \bigcup_{r_1 \in h_1(f)} h_2(r_1) \right).
\]
\end{definition}

\begin{definition}[Parallel]
Given DPs $d_1 = \langle F_1, R_1, h_1 \rangle$ and $d_2 = \langle
F_2, R_2, h_2 \rangle$, their \emph{parallel composition} is
$\langle F_1 \times F_2, R_1 \times R_2, h_1 \otimes h_2 \rangle$ where
\[
(h_1 \otimes h_2)(f_1, f_2) = \Min\!\left( h_1(f_1) \times h_2(f_2) \right).
\]
\end{definition}

\begin{definition}[Feedback]
Given a DP $d = \langle F \times X, R \times X, h \rangle$ with a
common factor $X$, the \emph{feedback composition} on $X$ is the DP
$d^{\dagger X} = \langle F, R, h^{\dagger X} \rangle$ where $h^{\dagger
X}(f)$ is the antichain obtained by closing the recursion on $X$.
\end{definition}

The recursion is solved by Kleene iteration in $\Antichain[R \times
X]$, seeded at $\{\bot\}$ and ascending until a fixed point is
reached.

\subsection{The Kleene Fixed-Point Algorithm}

\begin{theorem}[Censi 2015, Prop. 4]
\label{thm:kleene}
For an MCDP $d = \langle F \times X, R \times X, h \rangle$ closed
on $X$, define the operator $\Phi_f : \Antichain[R \times X] \to
\Antichain[R \times X]$ by
\[
\Phi_f(A) = \Min\!\left( \bigcup_{r \in A} h(f, r_X) \cap \uparrow\! r \right),
\]
where $r_X$ is the $X$-component of $r$ and $\uparrow\!r$ is the
upward closure. Then $\Phi_f$ is monotone, and the Kleene ascent
\[
A_0 = \{\bot\},\quad A_{k+1} = \Phi_f(A_k)
\]
converges to the least fixed point, which equals
$h^{\dagger X}(f)$ projected to $\Antichain[R]$.
\end{theorem}

In implementation, the iteration terminates when one of three
conditions holds: $A_{k+1} = A_k$ (fixed point), $A_{k+1} = \emptyset$
(infeasibility), or every point in $A_{k+1}$ has $r_X = \topp_X$
(loop axis saturates).

% ===========================================================================
\section{Installation and Project Layout}
\label{sec:install}
% ===========================================================================

\codename{} requires Python 3.9 or newer. The package has no
required runtime dependencies; \texttt{matplotlib} is needed only for
example~05.

\paragraph{Installation.} Clone the repository and install in editable
mode:
\begin{lstlisting}[language=bash]
git clone https://github.com/your-username/codesign-mcdp.git
cd codesign-mcdp
pip install -e .
\end{lstlisting}

\paragraph{Optional extras.}
\begin{itemize}[leftmargin=2em]
  \item \texttt{pip install -e ".[viz]"} adds matplotlib for the
        visualisation example.
  \item \texttt{pip install -e ".[nb]"} adds the Jupyter tooling needed
        to (re)generate the notebooks.
\end{itemize}

\paragraph{Repository layout.}
\begin{verbatim}
codesign-mcdp/
  codesign/          the Python package (9 modules)
  examples/          eight runnable example scripts
  notebooks/         the same examples as executed .ipynb notebooks
  tests/             smoke tests
  docs/              this manual plus rendered trace images
  build_notebooks.py script to regenerate the notebooks
  pyproject.toml     packaging metadata, PEP 621
  .github/workflows/ CI configuration
\end{verbatim}

% ===========================================================================
\section{Core Data Types}
\label{sec:posets}
% ===========================================================================

This section documents the data types provided by
\texttt{codesign.posets} and \texttt{codesign.antichains}.

\subsection{Posets}

The \texttt{Poset} abstract base class has the following interface:

\begin{lstlisting}
class Poset(ABC):
    name: str

    def leq(self, x, y) -> bool: ...      # is x <= y?
    def eq(self, x, y) -> bool: ...
    def lt(self, x, y) -> bool: ...
    def bottom(self): ...                  # least element, if any
    def top(self): ...                     # greatest element, if any
    def is_bottom(self, x) -> bool: ...
    def is_top(self, x) -> bool: ...
    def join(self, x, y): ...              # least upper bound
    def format(self, x) -> str: ...
\end{lstlisting}

Four concrete poset classes are provided.

\paragraph{\texttt{Reals(unit="")}.} The non-negative reals augmented
with $+\infty$. Bottom is $0.0$; top is \texttt{math.inf}. The optional
\texttt{unit} string is used for pretty-printing only.

\paragraph{\texttt{Naturals(unit="")}.} The non-negative integers
augmented with $+\infty$. Bottom is $0$; top is \texttt{math.inf}.

\paragraph{\texttt{NamedProduct(components: dict[str, Poset])}.}
The componentwise-ordered product of named factor posets. Elements
are Python dicts mapping component names to values. The bottom is the
dict of all bottoms; the top is the dict of all tops.
\texttt{NamedProduct} can be nested freely: a component poset may
itself be a \texttt{NamedProduct}.

\paragraph{\texttt{Discrete(elements, leq=None)}.} A finite poset with
explicit element set and a caller-supplied \texttt{leq} predicate.
Useful for catalog-like enumerations where the order is not the
default discrete order.

\subsection{Antichains}

The \texttt{Antichain} class wraps a set of poset elements together
with the poset they live in. The constructor normalises the input by
removing dominated points and deduplicating.

\begin{lstlisting}
class Antichain:
    poset: Poset
    points: tuple

    @classmethod
    def of_bottom(cls, poset) -> "Antichain": ...
    @classmethod
    def empty(cls, poset) -> "Antichain": ...
    @classmethod
    def singleton(cls, poset, point) -> "Antichain": ...
    @classmethod
    def from_set(cls, poset, points) -> "Antichain": ...

    def union_min(self, other) -> "Antichain": ...
    def filter_above(self, floor) -> "Antichain": ...
    def leq(self, other) -> bool: ...
    def eq(self, other) -> bool: ...
    def has_any_top(self) -> bool: ...
    def is_empty(self) -> bool: ...
\end{lstlisting}

The two key operations on antichains are \texttt{union\_min}, which
takes the union of two antichains and re-applies $\Min$, and
\texttt{filter\_above}, which keeps only points dominating a given
floor. Together they implement the body of the Kleene iteration in
Theorem~\ref{thm:kleene}.

% ===========================================================================
\section{Primitive Design Problem Types}
\label{sec:primitives}
% ===========================================================================

Six concrete subclasses of \texttt{DesignProblem} are provided, each
appropriate for a different style of relation between functionality
and resources.

\subsection{\texttt{AlgebraicDP}}

For relations where each resource is a closed-form monotone function
of the functionality.

\begin{lstlisting}
AlgebraicDP(
    F: NamedProduct,
    R: NamedProduct,
    equations: dict[str, Callable[[dict], float]],
    name: str = "algebraic",
)
\end{lstlisting}

The \texttt{equations} dictionary maps each resource name to a
callable taking a functionality dict and returning a scalar.
\texttt{h(f)} returns a singleton antichain.

\begin{example}
A battery sized by specific energy:
\begin{lstlisting}
battery = AlgebraicDP(
    F=NamedProduct({"capacity": Reals(unit="J")}),
    R=NamedProduct({"mass": Reals(unit="kg")}),
    equations={"mass": lambda f: f["capacity"] / 1.8e6},
)
\end{lstlisting}
\end{example}

\subsection{\texttt{FunctionDP}}

For relations where the user supplies \texttt{h} directly as a
function from functionality dicts to antichains. Use this when the
relation is multi-valued (a true Pareto front) or has branching logic.

\begin{lstlisting}
FunctionDP(
    F: NamedProduct,
    R: NamedProduct,
    h_fn: Callable[[dict], Antichain],
    name: str = "function",
)
\end{lstlisting}

\subsection{\texttt{CatalogDP}}

For selecting from a finite catalog of implementations (motors,
batteries, sensors).

\begin{lstlisting}
CatalogDP(
    F: NamedProduct,
    R: NamedProduct,
    catalog: list[CatalogEntry],
    name: str = "catalog",
)

CatalogEntry(
    name: str,
    provides: dict[str, Any],
    costs: dict[str, Any],
)
\end{lstlisting}

For a request $f$, $h(f)$ is the antichain of $\texttt{entry.costs}$
for every entry whose \texttt{provides} dominates $f$ in $F$. Entries
that are dominated by another feasible entry are pruned by $\Min$ in
$R$.

\subsection{\texttt{ConstraintDP}}

A feasibility predicate plus a scalar cost, lifted to $\Min$. Useful
when the resource-functionality relation is most natural to express
as \emph{is this choice good enough, and how much does it cost?}.

\begin{lstlisting}
ConstraintDP(
    F: NamedProduct,
    R: NamedProduct,
    sampler: Callable[[dict], Iterable[dict]],
    feasible: Callable[[dict, dict], bool],
    cost: Callable[[dict, dict], dict],
    name: str = "constraint",
)
\end{lstlisting}

The \texttt{sampler} yields candidate resource dicts for the given
functionality; \texttt{feasible} filters them; \texttt{cost} maps the
feasible ones into the resource poset, and the result is reduced by
$\Min$.

\subsection{\texttt{ODE\_DP}}

Derives a monotone resource relation from a differential equation
$\dot x = \text{rhs}(x, t, f)$. Two modes are supported.

\paragraph{Final value} integrates the ODE from $t = 0$ to
$t = t_{\text{end}}$ by explicit Euler and reports the extracted
resources of the final state.

\paragraph{Steady state} solves $\text{rhs}(x, \cdot, f) = 0$ for $x$
by Newton iteration and reports the extracted resources of the
steady-state value.

\begin{lstlisting}
ODE_DP(
    F: NamedProduct,
    R: NamedProduct,
    rhs: Callable,
    extract: Callable[[state], dict],
    mode: str,            # "final_value" or "steady_state"
    t_end: float = 1.0,
    n_steps: int = 100,
    x0_fn: Callable[[dict], Any] = ...,
    name: str = "ode",
)
\end{lstlisting}

\subsection{\texttt{UncertainDP}}

Brackets an unknown $h$ with a known lower and upper bound. The
optimistic mode (\texttt{lower}) yields a lower bound on the minimal
resources; the pessimistic mode (\texttt{upper}) yields an upper
bound. Section~VII of Censi (2015) develops the theory.

\begin{lstlisting}
UncertainDP(
    F: NamedProduct,
    R: NamedProduct,
    lower: DesignProblem,
    upper: DesignProblem,
    mode: str = "upper",
    name: str = "uncertain",
)
\end{lstlisting}

\texttt{with\_mode("lower" | "upper")} returns the bracket selected
for solving.

% ===========================================================================
\section{Composition Operators}
\label{sec:composition}
% ===========================================================================

\subsection{Series}

\begin{lstlisting}
series(dp1: DesignProblem, dp2: DesignProblem) -> Series
Series(dp1, dp2)
\end{lstlisting}

Requires \texttt{dp1.R == dp2.F} structurally (same component names
and posets). Implements the series definition: for each point of
$h_1(f)$ run $h_2$ and take the union $\Min$.

\subsection{Parallel}

\begin{lstlisting}
par(dp1: DesignProblem, dp2: DesignProblem) -> Parallel
Parallel(dp1, dp2)
\end{lstlisting}

Requires non-overlapping component names in $F_1, F_2$ and in
$R_1, R_2$. The output $F$ and $R$ are the concatenated
\texttt{NamedProduct}s. The antichain is the Cartesian product
followed by $\Min$.

\subsection{Loop}
\label{sec:loop}

\begin{lstlisting}
loop(inner: DesignProblem, axis: str) -> Loop
Loop(inner, axis="...")
\end{lstlisting}

Closes the feedback loop on a single named axis that must appear in
both \texttt{inner.F} and \texttt{inner.R}. The resulting DP has
\texttt{F = inner.F} minus the axis and \texttt{R = inner.R} minus
the axis. Calling \texttt{h(f)} triggers the Kleene fixed-point
iteration described in Theorem~\ref{thm:kleene}.

\begin{remark}[Multiple feedback loops]
\texttt{Loop} closes one axis at a time. To close several loops,
chain them: \texttt{loop(loop(inner, axis="x"), axis="y")}. The
\texttt{System} builder (Section~\ref{sec:system-builder}) closes all
feedback loops simultaneously over a single bundled axis, which is
usually more convenient.
\end{remark}

% ===========================================================================
\section{The Solver}
\label{sec:solver}
% ===========================================================================

\subsection{Top-level entry point}

\begin{lstlisting}
solve(
    dp: DesignProblem,
    functionality: dict | None,
    *,
    max_iter: int = 200,
    record_trace: bool = False,
) -> SolveResult
\end{lstlisting}

The result is a dataclass with the following fields:
\begin{lstlisting}
@dataclass
class SolveResult:
    antichain: Antichain      # the Pareto front (in dp.R)
    iterations: int           # Kleene steps taken (0 if no loop)
    converged: bool           # True if a fixed point was reached
    feasible: bool            # True if antichain has finite, nonempty points
    trace: list[Antichain]    # the sequence A_0, A_1, ... if recorded
\end{lstlisting}

\subsection{The Kleene iteration}

\begin{lstlisting}
kleene_loop(
    loop_dp: Loop,
    f_outer: dict,
    *,
    max_iter: int = 200,
    record_trace: bool = False,
    trace_out: list | None = None,
    info_out: dict | None = None,
) -> Antichain
\end{lstlisting}

The body of the iteration is direct from Theorem~\ref{thm:kleene}. The
implementation adds:
\begin{itemize}
  \item a \emph{divergence cap} of $10^{30}$ that converts an iterate
        exceeding the cap to $\topp$, so floating-point overflow
        becomes infeasibility rather than \texttt{inf}/\texttt{nan};
  \item \texttt{try}/\texttt{except} around the inner $h$ to catch
        \texttt{OverflowError}, \texttt{ValueError}, and
        \texttt{ZeroDivisionError}, again mapping to $\topp$;
  \item three termination conditions: fixed point reached, antichain
        empty (infeasible), or every point's loop axis equal to
        $\topp$ (provably infeasible).
\end{itemize}

\subsection{Scalar minimisation over an antichain}

\begin{lstlisting}
minimize_cost(
    result: SolveResult,
    cost_fn: Callable[[dict], float],
) -> dict | None
\end{lstlisting}

Returns the single resource bundle in \texttt{result.antichain} that
minimises \texttt{cost\_fn}, or \texttt{None} if the result is
infeasible. This is the scalarisation step that turns a Pareto front
into a single engineering choice.

\subsection{Observability: trace, verbose, callback, status}
\label{sec:solver-observability}

The solver can be made to report its progress in three independent
ways, and its termination reason is exposed through a structured
\texttt{status} field that is orthogonal to \texttt{feasible}.

\paragraph{The \texttt{status} field.} \texttt{SolveResult.status} is
one of three strings:
\begin{description}
\item[\texttt{"converged"}] The iteration reached a fixed point (or
provably hit $\top$) within \texttt{max\_iter} steps. The answer in
\texttt{result.antichain} is the algorithm's verdict; \texttt{feasible}
tells you whether that verdict is finite (a real design) or $\top$
(no design exists).
\item[\texttt{"max\_iter"}] The iteration was cut off at
\texttt{max\_iter}. The current antichain may or may not be near a
fixed point. Usually a sign to increase \texttt{max\_iter}, or to look
for an unintentional non-monotonicity in the model.
\item[\texttt{"diverged"}] At least one numeric component crossed the
divergence cap of $10^{30}$ before the iteration could settle, and the
solver stopped. Distinguishes pure numerical blow-up from clean
$\top$-infeasibility.
\end{description}
The previous Boolean \texttt{converged} field is preserved as a
backward-compatibility alias for \texttt{status == "converged"}.

\paragraph{Structured trace.} Passing \texttt{trace=True} to
\texttt{solve} populates \texttt{result.trace} with a list of
\texttt{TraceEntry} dataclasses, one per Kleene step (plus the seed at
iteration~0):
\begin{lstlisting}
@dataclass
class TraceEntry:
    iteration: int           # 0 = seed, then 1, 2, ...
    antichain: Antichain     # snapshot at this step
    n_points:  int           # convenience: len(antichain)
    delta:     float | None  # max absolute change (numeric)
                             # or 0/1 (discrete); None at iter 0
    elapsed_ms: float        # wall time for this step alone
\end{lstlisting}
With \texttt{trace=False} (the default) the field is \texttt{None}, so
the cost of unused tracing is zero.

\paragraph{Verbose printing.} The integer \texttt{verbose} controls
live output:
\begin{description}
\item[\texttt{verbose=0}] (default) silent.
\item[\texttt{verbose=1}] one summary line at the end, e.g.\
\texttt{[solve] converged: 17 iters, |A|=1, total=2.0ms, feasible=True}.
\item[\texttt{verbose=2}] one line per iteration, showing the delta,
the antichain size, and the step's wall time. Useful for watching
oscillating fixed points.
\end{description}

\paragraph{Iteration callback.} Passing
\texttt{on\_iteration=callable} registers a callback that receives
each \texttt{TraceEntry} as soon as it is produced. Useful for live
plots, custom logging, or interrupting an iteration that looks
pathological:
\begin{lstlisting}
def log_every_5(entry):
    if entry.iteration % 5 == 0:
        print(f"iter {entry.iteration}: delta={entry.delta:.3e}")

solve(dp, f, on_iteration=log_every_5)
\end{lstlisting}

% ===========================================================================
\section{Uncertainty}
\label{sec:uncertainty}
% ===========================================================================

Two kinds of parameter uncertainty are supported, declared on
\texttt{Module} instances and consumed by the same \texttt{solve}
entry point. Both are independent of the model's structure: a module
can have set-based uncertainty, stochastic uncertainty, both, or
neither.

\subsection{Declaration}

A \texttt{Module} carries optional attributes:
\begin{lstlisting}
class Battery(Module):
    F = {"capacity": Reals(unit="J")}
    R = {"mass":     Reals(unit="kg")}

    def __init__(self, specific_energy=1.8e6, efficiency=0.85):
        self.specific_energy = specific_energy
        self.efficiency = efficiency
        super().__init__()

    def h(self, f):
        return {"mass": f["capacity"]
                       / (self.specific_energy * self.efficiency)}

b = Battery()
b.uncertain_set  = Box(...)          # deterministic, set-based
b.uncertain_dist = Stochastic(...)   # statistical
\end{lstlisting}
The uncertainty solver discovers every such module by walking the
\verb|_codesign_modules| attribute that \texttt{System.build} attaches
to the returned DP. Before each solve, the module's named parameters
are reassigned; afterwards, their nominal values are restored. The
user's original \texttt{Battery()} instance is unchanged.

\subsection{Set-based uncertainty (deterministic)}

\paragraph{\texttt{Box}.} Axis-aligned interval product. Each
parameter is declared as \texttt{(lo, hi)} or \texttt{(lo, hi, direction)},
where \texttt{direction} is one of the four tokens
\texttt{"more\_is\_better"}, \texttt{"more\_is\_worse"},
\texttt{"less\_is\_better"}, \texttt{"less\_is\_worse"}. With all
directions declared, the worst case is a single corner read off in
constant time:
\begin{lstlisting}
b.uncertain_set = Box(
    specific_energy=(1.6e6, 2.0e6, "more_is_better"),
    efficiency=(0.80, 0.90, "more_is_better"),
)
# worst case = (specific_energy=1.6e6, efficiency=0.80)
\end{lstlisting}
When some directions are undeclared, all $2^n$ endpoint combinations
are evaluated and the worst is kept; cheap when $n$ is modest.

\paragraph{\texttt{Ellipsoid}.} An $n$-D ellipsoid
$(p - c)^\top \Sigma^{-1} (p - c) \leq 1$ in parameter space. The
covariance $\Sigma$ encodes both scale and correlation between
parameters, so an ellipsoid is a more honest model than a box when
parameters are believed to vary together:
\begin{lstlisting}
b.uncertain_set = Ellipsoid(
    center={"specific_energy": 1.8e6, "efficiency": 0.85},
    cov=[
        [1.0e10, -2.0e3],
        [-2.0e3,  2.5e-3],
    ],
    params=["specific_energy", "efficiency"],
    directions={
        "specific_energy": "more_is_better",
        "efficiency":      "more_is_better",
    },
)
\end{lstlisting}
With all directions declared, the worst case is the boundary point
$c + L\,u^\star$, where $L$ is the Cholesky factor of $\Sigma$ and
$u^\star$ is the unit vector in the direction of badness. Otherwise,
the boundary is sampled (controlled by \texttt{boundary\_samples}) and
the worst sample is returned.

\paragraph{2D conveniences.} \texttt{Disk(center, radius, ...)} and
\texttt{Circle(center, radius, ...)} reduce to the isotropic
ellipsoid $\Sigma = r^2 I$. For monotone systems with declared
directions, the two are equivalent: the worst case lies on the
boundary either way.

\subsection{Stochastic uncertainty}

\paragraph{\texttt{Stochastic}.} A joint distribution built from
named scipy-stats frozen marginals plus a copula:
\begin{lstlisting}
from scipy import stats

b.uncertain_dist = Stochastic(
    marginals={
        "specific_energy": stats.uniform(loc=1.6e6, scale=0.4e6),
        "efficiency":      stats.uniform(loc=0.80, scale=0.10),
    },
    copula=GaussianCopula(correlation=[[1.0, 0.4],
                                        [0.4, 1.0]]),
)
\end{lstlisting}
Sampling is by the standard copula--marginal trick: the copula
generates correlated uniforms in $[0,1]^d$; the marginal \texttt{ppf}s
map them to the named parameter axes.

\paragraph{Copulas.} \texttt{Independence()} is the default
(coordinatewise uniform). \texttt{GaussianCopula(correlation)} uses
the Gaussian copula with the given correlation matrix; samples are
drawn by Cholesky factorisation followed by $\Phi$ on each
coordinate. Subclassing \texttt{Copula} requires only one method,
\texttt{sample\_uniform(n, d, rng)}, so other copulas (e.g.\ Clayton,
Gumbel) drop in without further plumbing.

\subsection{Querying}

The same \texttt{solve} entry point accepts an \texttt{uncertainty}
list and returns a richer object:
\begin{lstlisting}
result = solve(
    dp, f,
    uncertainty=["worst_case", "mean", "p95", "cvar95", "samples"],
    n_samples=1000,
    rng_seed=42,
)

result.worst_case        # SolveResult at the worst point of the set
result.mean              # dict[r_port -> mean across MC samples]
result.p95               # dict[r_port -> 95th percentile]
result.cvar95            # dict[r_port -> CVaR at 95% level]
result.samples           # list[Antichain], one per MC sample
result.feasibility_rate  # fraction of feasible MC samples
\end{lstlisting}

The allowed labels are:
\begin{description}
\item[\texttt{"worst\_case"}] The deterministic worst over every
module's \texttt{uncertain\_set}. One canonical solve at the
worst-case parameter point.
\item[\texttt{"mean"}, \texttt{"p95"}, \texttt{"cvar95"}] Marginal
statistics over Monte Carlo samples, drawn from every module's
\texttt{uncertain\_dist}. The CVaR (conditional value at risk) at the
95\% level is the mean of the worst 5\% of samples.
\item[\texttt{"samples"}] The raw antichain per MC sample, for the
user to aggregate however they like.
\end{description}

For a single solve call you typically observe the ordering
\[
\text{nominal} < \text{mean} < \text{p95} < \text{CVaR95} < \text{worst\_case},
\]
which separates the four standard ways of reporting an answer "under
uncertainty." Notebooks 11 and 12 visualise this for the
example-7 drone.

% ===========================================================================
\section{The MCDP Builder}
\label{sec:mcdp-builder}
% ===========================================================================

The \texttt{MCDP} class in \texttt{codesign.mcdpl} provides an
MCDPL-style declarative syntax for a single self-contained design
problem.

\begin{lstlisting}
with MCDP("name") as m:
    m.provides("functionality_name", unit="...")    # outer F port
    m.requires("resource_name",     unit="...")    # outer R port

    m.constraint("resource_name", lambda f: <expr>)   # closed-form
    # or
    m.rule(lambda f: Antichain.from_set(...))         # multi-valued

    m.loop_on("axis_name")    # optional, may be repeated

dp = m.build()
\end{lstlisting}

The builder emits an \texttt{AlgebraicDP} or \texttt{FunctionDP}
internally, then wraps it in \texttt{Loop}s as requested by
\texttt{loop\_on}. Multiple \texttt{constraint} calls on the same
target are joined (max). Mostly useful when the model has a small
number of variables and reads naturally as a list of $\geq$
inequalities, mirroring the paper's \texttt{mcdp \{ ... \}} blocks.

% ===========================================================================
\section{The System Builder}
\label{sec:system-builder}
% ===========================================================================

For larger designs the \texttt{System} class in \texttt{codesign.system}
provides modular composition with named subsystems and algebraic
connection constraints. This is the recommended way to build
non-trivial co-design models. Two equivalent surface syntaxes are
available: the operator-overloaded form (recommended for new code) and
the legacy lambda form (still supported).

\subsection{Operator-overloaded syntax (recommended)}

\texttt{provides}, \texttt{requires}, and \texttt{add} each return a
port handle. Arithmetic operators on handles build expression trees
lazily; the \texttt{>=} operator registers a constraint with the parent
system.

\begin{lstlisting}
from codesign import Module, Reals, System, solve

class Battery(Module):
    F = {"capacity": Reals(unit="J")}
    R = {"mass":     Reals(unit="kg")}
    def h(self, f):
        return {"mass": f["capacity"] / 1.8e6}

class Actuator(Module):
    F = {"lift_force": Reals(unit="N")}
    R = {"power":      Reals(unit="W")}
    def h(self, f):
        return {"power": 10.0 * f["lift_force"] ** 2}

sys = System("drone")
endurance     = sys.provides("endurance",     unit="s")
extra_payload = sys.provides("extra_payload", unit="kg")
extra_power   = sys.provides("extra_power",   unit="W")
total_mass    = sys.requires("total_mass",    unit="kg")

battery  = sys.add("battery",  Battery())
actuator = sys.add("actuator", Actuator())

battery.capacity    >= (actuator.power + extra_power) * endurance
actuator.lift_force >= 9.81 * (battery.mass + extra_payload)
total_mass          >= battery.mass + extra_payload

drone = sys.build()
\end{lstlisting}

The four port kinds are:

\begin{itemize}[leftmargin=2em]
  \item \textbf{Outer F} (from \texttt{provides}): can appear on the RHS
        of constraint expressions. Cannot be a constraint target.
  \item \textbf{Outer R} (from \texttt{requires}): can be a constraint
        target. Cannot appear in expression RHS.
  \item \textbf{Module F} (\texttt{handle.f\_port}): can be a constraint
        target. Cannot appear in expression RHS (its value is determined
        by the constraints on it, not by the current iteration state).
  \item \textbf{Module R} (\texttt{handle.r\_port}): can appear in
        expression RHS. Cannot be a constraint target (its value is
        determined by the module's own \texttt{h}, not externally).
\end{itemize}

Misusing a port (e.g.\ trying to put a module F port on the RHS, or
constraining an outer F) raises immediately with a message naming the
port and explaining the rule. F and R port names within a single
subsystem must be disjoint.

The supported operators are \texttt{+, -, *, /, **,} unary \texttt{-},
along with the helper functions \texttt{codesign.sqrt}, \texttt{exp},
and \texttt{log} for expressions that need transcendental terms. For
demands that cannot be expressed algebraically, the legacy lambda form
remains available and can be mixed freely with the operator form (see
below).

\subsection{Legacy lambda syntax}

\begin{lstlisting}
sys.constrain("battery.capacity",
              lambda x: (x["actuator.power"] + x["extra_power"]) * x["endurance"])
sys.constrain("actuator.lift_force",
              lambda x: 9.81 * (x["battery.mass"] + x["extra_payload"]))
sys.constrain("total_mass",
              lambda x: x["battery.mass"] + x["extra_payload"])
\end{lstlisting}

The argument \texttt{x} is a dict carrying outer F values under their
bare names (\texttt{x["endurance"]}) and subsystem R ports under their
dotted names (\texttt{x["battery.mass"]}). Both syntaxes compile to the
same internal representation; the operator form is purely sugar.

\subsection{Class-based modules}

The \texttt{Module} base class lets a design problem be defined
declaratively:

\begin{lstlisting}
class Battery(Module):
    F = {"capacity": Reals(unit="J")}
    R = {"mass":     Reals(unit="kg")}

    def __init__(self, specific_energy=1.8e6):
        self.specific_energy = specific_energy
        super().__init__()

    def h(self, f):
        return {"mass": f["capacity"] / self.specific_energy}
\end{lstlisting}

A \texttt{Module} subclass declares its $F$ and $R$ ports as
class-level dicts and overrides \texttt{h(self, f)}. The constructor
wires everything into the underlying \texttt{DesignProblem} machinery.
Parameterised modules are written by overriding \texttt{\_\_init\_\_}
and calling \texttt{super().\_\_init\_\_()} after setting any
instance attributes. \texttt{h} may return a dict (treated as a
singleton antichain), a list of dicts (multi-valued antichain), or
an \texttt{Antichain} directly.

\subsection{Semantics}

The \texttt{build} method emits a \texttt{Loop} whose axis bundles
every subsystem's resource poset:
\[
\texttt{\_\_modules\_\_} =
  \prod_{m \in \mathrm{modules}} R_m.
\]
The inner DP performs one Kleene step:
\begin{enumerate}[leftmargin=2em]
  \item Read the current bundle estimate to populate
        $\mathtt{ctx[m.r]}$ for every subsystem $R$ port.
  \item For each subsystem $m$, evaluate the constraint demands on
        its $F$ ports to obtain $f_m$, then call $m.h(f_m)$ to obtain
        an antichain.
  \item Take the Cartesian product across subsystems; for each
        combination, evaluate the outer-$R$ constraint demands to
        produce a candidate point.
  \item Return the resulting antichain.
\end{enumerate}
The outer Kleene iteration converges over all subsystem $R$ ports
simultaneously.

\subsection{Validation}

\texttt{build} performs the following checks:
\begin{itemize}[leftmargin=2em]
  \item every subsystem $F$ port has at least one \texttt{constrain}
        target;
  \item every outer $R$ has at least one \texttt{constrain} target;
  \item every constraint target refers to a known module $F$ port or a
        declared outer $R$;
  \item module names do not clash with the reserved axis name
        \texttt{\_\_modules\_\_} and contain no \texttt{.}.
\end{itemize}
If any check fails, \texttt{build} raises \texttt{ValueError} with a
message identifying the missing or invalid item.

% ===========================================================================
\section{Worked Examples}
\label{sec:examples}
% ===========================================================================

The package ships twelve runnable examples under \texttt{examples/}
and the same twelve as executed notebooks under \texttt{notebooks/}.
Each is referenced below by its filename.

\subsection{Example 1: the drone (Fig. 48 of the paper)}

\texttt{01\_drone.py}. The canonical MCDP from the paper. A battery
powers an actuator that must lift the battery plus a payload, with
extra fixed payload and avionics power. The model is a single
\texttt{FunctionDP} closed by \texttt{Loop} on \texttt{battery\_mass}.
The example sweeps five mission profiles, three feasible and two
infeasible. Iteration counts range from 8 to 41 depending on how
close to the feasibility boundary the case is.

\subsection{Example 2: integer optimisation (Sec. VI-D)}

\texttt{02\_integer\_optimization.py}. The problem
\[
x + y \geq \lceil\sqrt{x}\rceil + \lceil\sqrt{y}\rceil + c
\]
over $\Nat \times \Nat$, parametrised by $c$. The inner DP enumerates
all valid splits of the deficit, producing a multi-point antichain at
each iteration. For $c = 4$ the iteration converges in six steps to
$M(4) = \{(0,7), (3,6), (4,4), (6,3), (7,0)\}$. The example notes
that the paper's claim $M(1) = \{(1,0), (0,1)\}$ contains a typo;
the correct answer is $\{(0,3), (3,0)\}$, which is what the solver
finds.

\subsection{Example 3: AUV seabed surveying (Sec. VIII)}

\texttt{03\_auv\_seabed.py}. An autonomous underwater vehicle sweeps
an area $A$ at velocity $v$ with sensor radius $r$. Cyclic constraints
couple time, energy, and cost. The example enumerates a small set of
$(v, r)$ tradeoffs per iteration, producing a three-point Pareto
front per scenario.

\subsection{Example 4: \texttt{UncertainDP} and \texttt{ODE\_DP}}

\texttt{04\_uncertain\_and\_ode.py}. Two demos of the advanced
primitives. First: a battery whose specific energy is known only to
lie in $[1.6, 2.0]$ MJ/kg; the lower and upper brackets give
optimistic and pessimistic masses. Second: a heater whose steady-state
input power is recovered from Newton's law of cooling via
\texttt{ODE\_DP}.

\subsection{Example 5: visualising the Kleene ascent}

\texttt{05\_visualize\_kleene.py}. Uses matplotlib to render the
sequence $A_0, A_1, \ldots$ for the Sec. VI-D problem, reproducing
the structure of Fig. 36 in the paper. PNGs for $c \in \{1, 4, 8\}$
are committed under \texttt{docs/images/}.

\subsection{Example 6: the drone in MCDPL syntax}

\texttt{06\_drone\_mcdpl\_syntax.py}. Rebuilds the drone with the
\texttt{MCDP} builder. The model reads almost identically to the
\texttt{mcdp \{ ... \}} block in Fig. 48 of the paper. Output is
identical to example~1.

\subsection{Example 7: the drone, modular}

\texttt{07\_drone\_modular.py}. Rebuilds the drone with the
\texttt{System} builder. Battery and actuator are defined as
\texttt{Module} subclasses with their own $F$ and $R$, then wired
together with three \texttt{>=} constraints in the operator-overloaded
syntax. Fixed-point values match example~1 exactly; iteration counts
are roughly $2\times$ because the bundle now has two coupled axes
(battery mass, actuator power) instead of one.

\subsection{Example 8: a modular vehicle with Pareto tradeoffs}

\texttt{08\_vehicle\_modular.py}. Three independent subsystems: a
\texttt{CatalogDP} motor (seven Pareto-incomparable entries), a
linear chassis, and a battery. Five \texttt{>=} constraints close the
loop on (motor, chassis, battery) resources. Two of the three mission
profiles yield two-point system Pareto fronts (different motor choices
trading total mass against total cost); the third is infeasible
because no motor in the catalog can supply enough torque once the
chassis grows to support the battery.

\subsection{Example 9: a robotic arm with non-trivial topology}

\texttt{09\_robotic\_arm.py}. Five subsystems (shoulder joint, elbow
joint, sensor, controller, battery) wired with eight constraints. The
connection graph is genuinely non-trivial: the controller's power
demand depends on both the sensor rate and the joint command rate, the
shoulder must support the elbow joint's mass at the end of its arm
(introducing a mechanical cycle through \texttt{elbow.mass}), and the
battery is sized by the integrated power of every electrical
subsystem. The Kleene iteration resolves the resulting cycles in four
steps for all three test mission profiles. This example is the cleanest
demonstration of where the operator-overloaded syntax pays off: each
constraint line states one physical relationship, and the cyclic
structure emerges from the constraint graph rather than from explicit
loop machinery.

\subsection{Example 10: watching the solver work}

\texttt{10\_solver\_trace.py}. The drone of example 7 instrumented to
show every observability feature of \texttt{solve}: silent mode, the
end-of-solve summary line (\texttt{verbose=1}), the per-iteration
progress feed (\texttt{verbose=2}), the structured trace
(\texttt{trace=True}), and a custom \texttt{on\_iteration} callback.
The same script runs the drone with three different \texttt{max\_iter}
caps and a deliberately under-resourced mission, demonstrating that
\texttt{status} takes the three distinct values \texttt{"converged"},
\texttt{"max\_iter"}, and \texttt{"diverged"} on these inputs. The
notebook version (\texttt{10\_solver\_trace.ipynb}) adds a
delta-vs-iteration plot showing the characteristic oscillating decay
of the Kleene iteration to machine precision.

\subsection{Example 11: set-based deterministic uncertainty}

\texttt{11\_uncertain\_drone.py}. The drone's battery has two
internal parameters (specific energy, efficiency) declared with
"more is better" directions. First a \texttt{Box} with independent
ranges, then a tilted \texttt{Ellipsoid} of the same widths. The
worst-case mass under the box is +0.0158\,kg above nominal; under the
ellipsoid it is +0.0050\,kg. The discrepancy is the cost of admitting
the implausible corner where both parameters are simultaneously at
their worst; the ellipsoid model rejects it explicitly.

\subsection{Example 12: stochastic uncertainty with a Gaussian copula}

\texttt{12\_stochastic\_drone.py}. The same drone with two uniform
marginals tied by a Gaussian copula at correlation~0.4. A single
\texttt{solve} call requests all five summaries:
\texttt{worst\_case}, \texttt{mean}, \texttt{p95}, \texttt{cvar95},
\texttt{samples}. Over a thousand Monte Carlo samples, the five
values come out as nominal $<$ mean $<$ p95 $<$ CVaR95 $<$
worst\_case, the canonical ordering. The notebook adds a matplotlib
histogram of the MC distribution with vertical markers at each
summary.

% ===========================================================================
\section{Modelling Guidelines and Pitfalls}
\label{sec:guidelines}
% ===========================================================================

A few patterns that come up repeatedly when authoring MCDPs.

\paragraph{Expose loop variables you want to inspect.} When using the
operator API directly, \texttt{Loop} projects its axis out of the
outer $R$. To inspect the converged loop value, include it in the
inner $R$ under a \emph{different name}: for example, both
\texttt{battery\_mass} as the loop axis and \texttt{report\_mass}
mirrored for outer-$R$ visibility. The \texttt{System} builder
handles this transparently.

\paragraph{Cap physical maxima.} When a design variable has a
physical ceiling ($v_{\max}$, $r_{\max}$), make $h$ return a
$\topp$-valued antichain whenever the iterate exceeds it. The Kleene
ascent will then converge to infeasibility rather than oscillating
or diverging.

\paragraph{Use \texttt{FunctionDP} or \texttt{CatalogDP} for breadth.}
\texttt{AlgebraicDP} always returns a singleton antichain. For a
genuine Pareto front, use \texttt{FunctionDP} (enumerate the
tradeoffs explicitly) or \texttt{CatalogDP} (provide multiple
incomparable entries).

\paragraph{Scalarisation is downstream.} The MCDP solver returns the
Pareto front. Choosing which point to ship is a separate concern:
apply \texttt{minimize\_cost} with a scalar objective, or inspect the
front and pick by hand.

\paragraph{Catalog Cartesian product can blow up.} When several
subsystems return multi-valued antichains, the inner $h$ of the
\texttt{System} takes their Cartesian product. For typical models
(one or two catalog subsystems among many algebraic ones), this is
not a bottleneck; but if many subsystems are catalog-driven, the
inner antichain may grow large between $\Min$ prunings.

\paragraph{Numerical precision in feedback loops.}
\texttt{Antichain.eq} compares points exactly. For contractive
floating-point feedback, the iteration may oscillate by an ulp near
the fixed point and consume many iterations before converging. The
divergence cap and \texttt{max\_iter} guard against runaway, but
choosing \texttt{max\_iter} generously (200 to 500) is wise for
nontrivial models.

% ===========================================================================
\section{Limitations and Future Work}
\label{sec:limitations}
% ===========================================================================

Version 0.1.0 covers the algorithmic core: posets, antichains, six
primitive DP types, three composition operators, the Kleene solver,
and the two builders. The following are intentional omissions, all
tractable extensions:

\begin{itemize}[leftmargin=2em]
  \item \textbf{MCDPL surface syntax}. The paper's
        \texttt{mcdp \{ ... \}} text format is not parsed. The
        \texttt{MCDP} builder gives the same shape in Python.
  \item \textbf{Non-finitely-representable antichains}. The library
        works on finite, exactly-represented antichains. For
        relations whose Pareto front is continuous, the only support
        is the bracket pattern of \texttt{UncertainDP}. Approximation
        kernels (Sec.~VII of Censi 2015) are not yet implemented.
  \item \textbf{Visualisation of the design graph}. Example~5
        visualises the Kleene ascent over $\Nat \times \Nat$, but
        there is no general renderer for the operator tree or for
        the subsystem network of a \texttt{System}.
  \item \textbf{Symbolic/automatic differentiation of $h$}. All
        relations are evaluated numerically.
  \item \textbf{Distributed or parallel solving}. The Kleene
        iteration is sequential.
\end{itemize}

Contributions, particularly on these items, are welcome.

% ===========================================================================
\section{API Reference Summary}
\label{sec:api}
% ===========================================================================

The public API exported by \texttt{import codesign} is summarised
below.

\paragraph{Posets.}
\texttt{Reals}, \texttt{Naturals}, \texttt{NamedProduct},
\texttt{Discrete}, abstract base \texttt{Poset}.

\paragraph{Antichains.}
\texttt{Antichain}.

\paragraph{Design problems.}
\texttt{DesignProblem} (abstract base),
\texttt{AlgebraicDP},
\texttt{FunctionDP},
\texttt{CatalogDP}, \texttt{CatalogEntry},
\texttt{ConstraintDP},
\texttt{ODE\_DP},
\texttt{UncertainDP}.

\paragraph{Composition.}
\texttt{Series}, \texttt{Parallel}, \texttt{Loop} (classes);
\texttt{series}, \texttt{par}, \texttt{loop} (function aliases).

\paragraph{Primitives.}
\texttt{adder}, \texttt{multiplier}, \texttt{scale}, \texttt{constant},
\texttt{identity}.

\paragraph{Solver.}
\texttt{solve}, \texttt{kleene\_loop}, \texttt{minimize\_cost},
\texttt{SolveResult}, \texttt{TraceEntry}.

\paragraph{Builders.}
\texttt{MCDP}, \texttt{System}.

\paragraph{Class-based modules.}
\texttt{Module} (declarative \texttt{DesignProblem} base class).

\paragraph{Constraint DSL.}
\texttt{Port}, \texttt{Expr}, \texttt{ModuleHandle} (returned by
\texttt{System.add}), and the helper functions \texttt{sqrt},
\texttt{exp}, \texttt{log} for use inside constraint expressions.

\paragraph{Uncertainty.}
Set-based: \texttt{UncertaintySet} (abstract), \texttt{Box},
\texttt{Ellipsoid}, \texttt{Disk}, \texttt{Circle}.
Stochastic: \texttt{Stochastic}, plus copulas
\texttt{Independence} and \texttt{GaussianCopula}.
Result: \texttt{UncertaintyResult}.
Entry point: \texttt{solve\_with\_uncertainty} (the same machinery is
also reachable through \texttt{solve(dp, f, uncertainty=[...])}).

\paragraph{Module-level.}
\texttt{\_\_version\_\_} (the string \texttt{"0.1.0"}).

% ===========================================================================
\section*{References}
% ===========================================================================
\addcontentsline{toc}{section}{References}

\begin{enumerate}[label={[\arabic*]},leftmargin=2em]
  \item Andrea Censi. \emph{A Mathematical Theory of Co-Design}.
        arXiv:1512.08055, 2015.
        \url{https://arxiv.org/abs/1512.08055}
  \item B.\,A. Davey and H.\,A. Priestley. \emph{Introduction to
        Lattices and Order}. Cambridge University Press, 2nd ed.,
        2002.
  \item The MCDPL software distribution.
        \url{https://co-design.science/software/}
\end{enumerate}

\end{document}
